% ===================================================================
% ENGLISH ABSTRACT
% ===================================================================

\chapter{Abstract}

Magnetic Resonance Imaging (MRI) has become an indispensable diagnostic tool for clinicians and a powerful investigative technique for researchers. Even after more than five decades since its inception, the field of MRI continues to undergo substantial and consistent advances in technology, methodology, and applications. Using nuclei ranging from the abundant hydrogen (¹H) to less abundant X-nuclei such as deuterium (²H), cabon (13C) sodium (²³Na), and phosphorus (³¹P), MRI enables the measurement of structures from macro to micro scale and the assessment of a wide range of physiological parameters, including blood oxygenation, metabolism, perfusion, pH, and temperature, with varying degrees of success. At present, a significant research is focused on ultra high field MRI ($\geq 7\,T$). The greatly improved sensitivity and distinctive contrast mechanisms at such high fields have expanded the capabilities of structural and functional imaging, particularly of the brain. This work primarily focuses on the methodological development of a rapid acquisition scheme called balanced Steady State Free Precession (bSSFP) for investigating brain function. bSSFP is known for its high signal to noise ratio (SNR) efficiency and its characteristic T2 over T1 contrast, especially in fluid rich tissues, which has traditionally made it useful in cardiac imaging at lower field strengths.

In the first part of this thesis was focused on improving acquisition efficiency of bSSFP-based functional MRI (fMRI) using fast and efficient non-Cartesian trajectories with parallel imaging and investigate their effects on the measured blood oxygen level dependent (BOLD) contrast. The aim is to have a distortion-free fMRI method with improved spatial localization of neuronal activity for \gls{UHF} application  as an alternative to the popular gradient echo EPI. The micro-vascular sensitivity of bSSFP, similar to that of spin echo sequences, is expected to provide better localization of the BOLD signal to sites of neuronal activation. Moreover, the inherently lower radiofrequency (RF) power deposition of bSSFP is advantageous under the restrictive specific absorption rate (SAR) limits at high fields. However, achieving high quality images at such resolutions demands stringent performance from the imaging hardware. Therfore,  system imperfections including eddy current effects and spatial $B_0$ inhomogeneity were included in the reconstuction model to improve image quality. The flexibility of the developed acquisition and reconstruction method was tested

Secondly, bSSFP was employed to improve the sensitivity deuterium-labeled Gluocose 



short T1 and large T2/T1 ratio yield a significant SNR boost
large B0
healthy human brain 

In summary, two aspect of bSSFP i.e spin-echo behavior to improve speficity of fMRI and high signal efficiency of large T2/T1 species to improve human DMI was explored to investigate brain function at an ultra high field of 9.4 T.

%%% Local Variables:
%%% mode: latex
%%% TeX-master: "../thesis"
%%% End:
